%------------------------------------------------------------------------------%
% ABSTRACT HARMONIC ANALYSIS                                                   %
%------------------------------------------------------------------------------%
% File  : Abstract_Harmonic_Analysis.tex                                       %
% Author: Klaus Hoeflinger, Beethovenstrasse 11, D-73117 Wangen                %
% Email : khoeflinger@t-online.de                                              %
%------------------------------------------------------------------------------%

%------------------------------------------------------------------------------%
%                       DOCUMENT CLASS and INCLUDE FILES                       %
%------------------------------------------------------------------------------%
\documentclass[11pt,twoside]{article}
\usepackage{geometry}
\geometry{paperheight=241mm,
          paperwidth=152mm,
          top=22mm,
          bottom=17mm,
          left=20mm,
          right=20mm,
          textwidth=112mm,
          textheight=202mm,
          headheight=10mm,
          headsep=3mm,
          footskip=9mm,
          includehead,
          includefoot,
          marginparsep=0mm,
          marginparwidth=0mm}

\renewcommand{\baselinestretch}{0.945}\normalsize

\AtBeginDocument{\setlength\abovedisplayskip{2mm}}
\AtBeginDocument{\setlength\belowdisplayskip{2mm}}

%------------------------------------------------------------------------------%
% define title formats                                                         %
%------------------------------------------------------------------------------%
\usepackage{titlesec}

\renewcommand{\thesection}{\arabic{section}}
\renewcommand{\sfdefault}{FiraSans}
\renewcommand{\ttdefault}{pcr}
\renewcommand{\rmdefault}{antiqua}

\titleformat{\part}[display]
{\normalfont \large \rmfamily \bfseries \centering}
{\small{\thepart.}}{2pt}{}

\titleformat{\section}[runin]
{\normalfont \rmfamily \bfseries}
{\thesection}{1pt}{.\;}[.]

%\titleformat{\section}
%{\normalfont \bfseries \small \filcenter}
%{\thesection.\;}{0mm}{}

\titleformat{\subsection}[runin]
{\normalfont \itshape}
{}{0mm}{}

\titleformat{\subsubsection}[runin]
{\normalfont \itshape}
{}{}{}{}

\titleformat{\paragraph}[runin]
{\normalfont}
{}{}{}{}

\titleformat{\subparagraph}[runin]
{\normalfont}
{}{}{}{}

\titlespacing*{\part}          { 0pt}{ 0pt}{ 8pt}
\titlespacing*{\section}       { 0pt}{ 7pt}{ 4pt}
\titlespacing*{\subsection}    { 0pt}{14pt}{ 6pt}
\titlespacing*{\subsubsection} { 0pt}{ 6pt}{12pt}
\titlespacing*{\paragraph}     { 0pt}{ 6pt}{12pt}
\titlespacing*{\subparagraph}  { 0pt}{ 0pt}{12pt}

%------------------------------------------------------------------------------%
% define enumerate parameters                                                  %
%------------------------------------------------------------------------------%
\usepackage{enumitem}
\setlength{\parindent}{3pt}
\setlength{\parskip}  {0pt}

%\setlist{noitemsep}

\setlist[enumerate]{
         leftmargin=12pt,
         rightmargin=0pt,
         itemsep=1pt,
         itemindent=3pt,
         topsep=3pt,
         labelsep=4pt,
         partopsep=0pt,
         labelwidth=0pt,
         listparindent=0pt,
         parsep=0pt}

\setlist[itemize]{
         leftmargin=12pt,
         rightmargin=0pt,
         itemsep=0pt,
         itemindent=1pt,
         topsep=1pt,
         labelsep=4pt,
         partopsep=1pt,
         labelwidth=0pt,
         listparindent=0pt,
         parsep=0pt}

%------------------------------------------------------------------------------%
% define packages                                                              %
%------------------------------------------------------------------------------%
\usepackage{upref}
\usepackage{amssymb}
\usepackage{slantsc}
\usepackage{amsmath}
\usepackage{amsthm}
\usepackage{thmtools}
\usepackage{amscd}
\usepackage{verbatim}
\usepackage{latexsym}
\usepackage{multicol}
\usepackage{graphicx}
\usepackage{setspace}
\usepackage[ngerman,english]{babel}
\usepackage[protrusion=true,expansion=true]{microtype}
\usepackage{tikz}
\usetikzlibrary{arrows,arrows.meta,decorations.markings}
\usepackage{mathrsfs}
\usepackage{amsfonts}
\usepackage{datetime2}
\usepackage{textgreek}
\usepackage{hyperref}
\usepackage{pifont}
\usepackage{relsize}
%\usepackage{biblatex}
\relsize{-0.05}
\usepackage[skip=0mm]{parskip}
\usetikzlibrary{decorations.pathmorphing}

\usepackage{mlmodern}
\usepackage[T5,T1]{fontenc}
\usepackage{ragged2e}

%\usepackage{mathabx}
%\usepackage{times}

%\usepackage{fontspec}
%\setmainfont{Nimbus Roman No9 L}
%\usepackage[T1]{fontenc}

%------------------------------------------------------------------------------%
% define page layout                                                           %
%------------------------------------------------------------------------------%
\usepackage{fancyhdr}
\pagestyle{fancy}
\setlength{\headheight}{30.0pt}
\renewcommand{\headrulewidth}{0pt}

\AtBeginDocument{
  \addtolength\abovedisplayskip{-0.0\baselineskip}
  \addtolength\belowdisplayskip{-0.0\baselineskip}
}

\fancyhead{}
\fancyfoot{}
\addtolength{\headwidth}{0.02\marginparwidth}

\makeatletter
\let\ps@plain\ps@fancy
\makeatother

\renewcommand\sectionmark[1]
{
 \def\sectionname{#1}
 \markright{\thesection #1}
}

\makeatletter
\@addtoreset{section}{part}
\makeatother

%------------------------------------------------------------------------------%
% define footnote without number                                               %
%------------------------------------------------------------------------------%
\newcommand{\myfootnote}[1]{%
\renewcommand{\thefootnote}{}%
\footnotetext{#1}%
\renewcommand{\thefootnote}{\arabic{footnote}}%
}

%------------------------------------------------------------------------------%
% define numbering in equations                                                %
%------------------------------------------------------------------------------%
\numberwithin{equation}{section}

%------------------------------------------------------------------------------%
% define newtheorem                                                            %
%------------------------------------------------------------------------------%
\declaretheoremstyle[
headfont=\scshape, 
headindent = 0mm, %\parindent,
%postheadhook = {\hspace{0mm}\newline},
%postheadspace = \newline, 
spaceabove = 3mm, 
spacebelow = 3mm,
numberwithin=section,
name=Theorem]{khMATH}
\declaretheorem[style=khMATH]{theorem}

\declaretheoremstyle[
headfont=\bfseries, 
headindent = 0mm, %\parindent,
%postheadhook = {\hspace{0mm}\newline},
%postheadspace = \newline, 
spaceabove = 3mm, 
spacebelow = 3mm,
numberwithin=part,
name=Theorem]{khMATH1}
\declaretheorem[style=khMATH1]{theorem1}

\declaretheoremstyle[
headfont=\scshape, 
headindent = 0mm, %\parindent,
%postheadhook = {\hspace{0mm}\newline},
%postheadspace = \newline, 
spaceabove = 3mm, 
spacebelow = 3mm,
numberwithin=section,
name=Corollary]{khMATH1}
\declaretheorem[style=khMATH1]{corollary}

%            name         counter     label              numeration-mode
\theoremstyle{plain}
\newtheorem*{introduction}            {\sc Introduction}
\newtheorem {standard}                {}                 [section]
\newtheorem*{definition}              {\sc Definition}
\newtheorem{satz1}                    {\sc Satz}
\newtheorem{lemma}                    {\sc Lemma}
\newtheorem*{proposition}             {\sc Proposition}
%\newtheorem{corollary}               {\sc Corollary}
\newtheorem*{corollar}                {\sc Korollar}
\newtheorem*{remark}                  {\sc Remark}
\newtheorem*{remarks}                 {\sc Remarks}
\newtheorem*{example}                 {Example}
\newtheorem*{examples}                {Examples}
\newtheorem*{theo}                    {\sc Theorem}

%------------------------------------------------------------------------------%
% define tabular                                                               %
%------------------------------------------------------------------------------%
\usepackage{tabularx}
\newcolumntype{L}[1]{>{\raggedright\arraybackslash}p{#1}}
\newcolumntype{C}[1]{>{\centering\arraybackslash}  p{#1}} 
\newcolumntype{R}[1]{>{\raggedleft\arraybackslash} p{#1}} 

%------------------------------------------------------------------------------%
% define \sh, \zB                                                              %
%------------------------------------------------------------------------------%
\def\dsh{{d.\,h.}}
\def\ie{{i.\,e.}}
\def\zB{z.\,B.}
\renewcommand{\abstractname}{ }

%------------------------------------------------------------------------------%
% change default spacing for cases                                             %
%------------------------------------------------------------------------------%
\makeatletter
\renewcommand*\env@cases[1][1.6]{%
  \let\@ifnextchar\new@ifnextchar
  \left\lbrace
  \def\arraystretch{#1}%
  \array{@{}l@{\quad}l@{}}%
}
\makeatother

\newenvironment{rcases}
  {\left.\begin{aligned}}
  {\end{aligned}\right\rbrace}

%------------------------------------------------------------------------------%
% change default for \predisplaypenalty                                        %
%------------------------------------------------------------------------------%
\predisplaypenalty=990
\allowdisplaybreaks \numberwithin{equation}{section}

%------------------------------------------------------------------------------%
% general notations                                                            %
%------------------------------------------------------------------------------%
\def\*{\star}
\def\={\,=\,}
\def\+{\,+\,}
\def\xsubset{{\,\subset\,}}
\def\xtimes{{\,\times\,}}
\def\xeof{{\,\in\,}}
\def\eq{{\,=\,}}
\def\xto{{\,\to\,}}
\def\eqdef{\,:=\,}
\def\:{{\,:\,}}
\def\ele{{\,\in\,}}
\def\exists{\mbox{ exists }}
\def\and{\mbox{ and }}
\def\for{\mbox{ for }}
\def\fa{\mbox{ for all }}
\def\fe{\mbox{ for each }}
\def\qfa{\quad \mbox{ for all }}
\def\df{\,\dotfill}
\def\iff{if and only if }
\def\wrt{with respect to}

\makeatletter
\newcommand \Df {\leavevmode \cleaders \hb@xt@ .70em{\hss .\hss }\hfill \kern \z@}
\makeatother

\def\sql{\mathfrak{a}}               % sesquilinear form a
\def\DF{B}                           % Dirichlet form a
%------------------------------------------------------------------------------%
% number sets and special elements                                             %
%------------------------------------------------------------------------------%
\def\P{\mathbb{H}}                   % half plane of $\C$
\def\J{I}                            % interval I
\def\N{{\mathbb{N}}}                 % unsigned integers
\def\Q{{\mathbb{Q}}}                 % rational integers
\def\Z{{\mathbb{Z}}}                 % integers
\def\R{{\mathbb{R}}}                 % real numbers
\def\Rn{{\mathbb{R}}^n}              % real numbers in Rn
\def\Rd{{\mathbb{R}}^d}              % real numbers in Rn
\def\C{{\mathbb{C}}}                 % complex numbers
\def\F{{\mathbb{K}}}                 % arbitrary field
\def\Torus{{\mathbb{T}}}             % complex circle line
\def\T{\tau}                         % upper bound of interval (0,t)
\def\sign#1{\operatorname{sign}(#1)} % sign of a number

%------------------------------------------------------------------------------%
% set theory                                                                   %
%------------------------------------------------------------------------------%
\def\G{G}                            % domain $G$
\def\clG{\overline{\G}}              % closure of $G$
\def\bndG{\partial \G}               % boundary of $G$
\def\setA{\mathcal{A}}               % set A
\def\setB{\mathcal{B}}               % set B
\def\setC{\mathcal{C}}               % set C
\def\setM{M}                         % set M
\def\setN{N}                         % set N
\def\setS{\mathcal{S}}               % set S
\def\famF{\mathfrak{F}}              % family of sets
\def\indI{\mathcal{I}}               % index set I
\def\pot#1{\mathfrak{P}(#1)}         % power set

\def\relR{\mathbf{R}}                % relation R
\def\mapf{f}                         % mapping f
\def\mapg{g}                         % mapping g

\def\set#1#2{\{ \, #1 \,: \, #2 \, \}}
\def\tensor#1#2{#1 \otimes #2}
\def\Tens{\oplus}

\def\cal#1{{\mathcal{#1}}}
\def\aut#1{\operatorname{Aut}(#1)}

\def\cross#1#2{#1 {\,\times\,} #2}
\def\multicross#1#2{#1 \times  \ldots \times  #2}
\def\multitens#1#2 {#1 \otimes \ldots \otimes #2}

%------------------------------------------------------------------------------%
% group theory                                                                 %
%------------------------------------------------------------------------------%
\def\1{{\mathsf{1}}}                % unit element of a monoid
\def\0{{\mathsf{0}}}                % unit element of an Abelian monoid
\def\kernel#1{\mathfrak{N}(#1)}     % kernel of a homomorphism
\def\cokernel#1{\mathfrak{C}{(#1)}} % cokernel of a homomorphism
\def\Hom#1{\operatorname{Hom}(#1)}

%------------------------------------------------------------------------------%
% linear algebra                                                               %
%------------------------------------------------------------------------------%
\def\vecV{\mathit{V}}               % vector space V
\def\vecH{\mathit{H}}               % vector space H
\def\vecW{\mathit{W}}               % vector space W
\def\vecU{\mathit{U}}               % subspace U
\def\convC{\mathcal{C}}             % convex set C
\def\basH{\mathfrak{B}}             % Hamel base B

\def\LO#1{\mathscr{L}(#1) }         % algebra of linear transformations

\def\scalar#1#2{ \langle #1,#2 \rangle }
\def\trace#1{\operatorname{tr}(#1)}
\def\dim#1{{\operatorname{dim} \, #1 }}
\def\codim#1{{\operatorname{codim} \, #1 }}
\def\dim{\operatorname{dim}}
\def\codim{\operatorname{codim}}
\def\ind{\operatorname{ind}}

\def\genG{\cal{G}}

\def\algA{\cal{A}}
\def\algB{\cal{B}}
\def\algU{\cal{U}}

\def\idI{\cal{I}}

%------------------------------------------------------------------------------%
% topology                                                                     %
%------------------------------------------------------------------------------%
\def\compK{{K}}          % compact set C
\def\openO{\mathcal{O}}  % open set O
\def\openU{\mathcal{U}}  % open set U
\def\U{\mathcal{U}}      % neighbourhood U
\def\V{\mathcal{V}}      % neighbourhood V
\def\d{\mathit{d}}       % metric function d

\def\interior#1{\mathring{#1}}
\def\closure#1{{\overline{#1}}}

\def\topT{\mathcal{T}}   % topology T
\def\topX{X}             % topological space X
\def\topY{Y}             % topological space Y
\def\topZ{Z}             % topological space Z
\def\metrS{M}            % metric space M

\def\grpG{G}             % topological group G
\def\grpA{A}             % topological group A
\def\grpB{B}             % topological group B
\def\grpC{C}             % topological group C
\def\grpH{H}             % topological group H
\def\grpU{U}             % topological group U
\def\grpV{V}             % topological group V
\def\topG{G}             % topological group G
\def\topH{H}             % topological group H
\def\topU{U}             % topological subgroup U

\def\subS{\cal{S}}
\def\riR{R}              % ring R

%------------------------------------------------------------------------------%
% calculus                                                                     %
%------------------------------------------------------------------------------%
\def\supp#1{\operatorname{supp}(#1)}
\def\singsupp#1{\operatorname{sing \, supp}(#1)}

\def\re{\operatorname{Re}}
\def\im{\operatorname{Im}}
\def\grad{\operatorname{grad}}

%\def\abs#1 {\left|#1\right|}
\def\abs#1 {\lvert #1 \rvert}
%\def\abs#1 {\left|\,#1\,\right|}

%------------------------------------------------------------------------------%
% measure theory                                                               %
%------------------------------------------------------------------------------%
\def\salgA{\Sigma}               % sigma-algebra S
\def\salgB{\cal{B}}              % sigma-algebra B
\def\Borel#1{\cal{B}_{#1}}       % Borel-algebra 
\def\LB{\lambda}                 % Lebesgue-Borel measure

%------------------------------------------------------------------------------%
% function spaces                                                              %
%------------------------------------------------------------------------------%
\def\Lp{L^p(\G)}                 % spaces L_p
\def\Lq{L^q(\G)}                 % spaces L_q
\def\Lh{L^2(\G)}                 % spaces L_2
\def\Li{L^{\infty}(\G)}          % spaces L_infty
\def\Lc{L_{loc}^1(\G)}           % spaces L_1^{loc}
\def\Ck{C^k(\G)}                 % spaces C_k

\def\BV#1{BV(#1) }               % set of functions of bounded variation
\def\AC#1{AC(#1) }               % set of absolutely continuous functions

\def\MP#1{\mathscr{M}^+(#1) }    % set of positive measures
\def\MS#1{\mathscr{M}^-(#1) }    % vector space of finite signed measures
\def\MC#1{\mathscr{M}(#1) }      % vector space of complex measures
\def\MCR#1{\mathscr{M}_{r}(#1) } % vector space of regular complex measures

\def\SobW{W^{k,p}(\G)}                % Sobolev space W
\def\SobO{\mathring{W}^{k,p}(\G)}     % Sobolev space W
%\def\WN#1#2{\mathring{W}^{#1}(#2) }   % Sobolev space W
\def\WN#1#2{W_0^{#1}(#2) }   % Sobolev space W
%\def\HN#1#2{\mathring{H}^{#1}(#2) }   % Sobolev space H
\def\HN#1#2{H_0^{#1}(#2) }   % Sobolev space H
%------------------------------------------------------------------------------%
% functional analysis                                                          %
%------------------------------------------------------------------------------%
\def\snorm#1{\lVert #1 \rVert}
\newcommand{\norm}[1]{\left\lVert #1 \right\rVert}

\def\topV{V}                       % topological vector space V
\def\topW{W}                       % topological vector space W
\def\normS{X}                      % normed space

\def\banX{\mathit{X}}              % Banach space X
\def\banY{\mathit{Y}}              % Banach space Y
\def\banZ{\mathit{Z}}              % Banach space Z
\def\dbanX{\mathit{X^*}}           % dual space of Banach space X
\def\dbanY{\mathit{Y^*}}           % dual space of Banach space Y
\def\dbanZ{\mathit{Z^*}}           % dual space of Banach space Z

\def\banA{\mathit{A}}              % Banach algebra A
\def\banB{\mathit{B}}              % Banach algebra B
\def\banC{\mathit{C}}              % C*-algebra C


\def\domain#1{\mathfrak{D}(#1)}    % domain
\def\range#1{\mathfrak{R}(#1)}     % range

\def\isoJ{\mathfrak{J}}            % isometry J
\def\repA{{\cal{A}_{\sql}}}        % representation operator A
\def\linS{{S}}                     % linear operator S
\def\linG{{G}}                     % linear operator G
\def\linL{{L}}                     % linear operator L
\def\linT{{T}}                     % linear operator A
\def\linA{\mathit{A}}              % linear operator A
\def\linB{{B}}                     % linear operator B
\def\linM{{M}}                     % linear operator M
\def\linU{{U}}                     % linear operator U
\def\linI{{I}}                     % linear operator I
\def\A{\cal{A}}                    % linear differential operator A
\def\BDO{\cal{B}}                  % linear boundary differential operator B
\def\linU{{U}}                     % unitary operator U
\def\linP{{P}}                     % projection P
\def\compA{k}                      % compact operator k
\def\linFR{f}                      % finite rank operator f
\def\fredA{A}                      % Fredholm operator F

\def\HAM{\cal{H}}                  % Hamilton operator
\def\PA{\partial^{\alpha}}         % elementary differential operator
\def\DO{\cal{A}}                   % linear differential operator
\def\BC{\cal{B}}                   % boundary operator
\def\SLO{\cal{L}_{p,q}}            % Sturm-Liouville operator
\def\MO{\cal{L}_{M}}               % momentum operator

\def\HS#1{S(#1)           }        % algebra of Hilbert-Schmidt operators
\def\FR#1{F(#1)           }        % algebra of finite rank operators
\def\TC#1{N(#1)           }        % algebra of trace class operators
\def\BU#1{\mathscr{U}(#1) }        % algebra of unitary linear operators
\def\BO#1{\mathscr{L}(#1) }        % algebra of bounded linear operators
\def\CO#1{\mathscr{C}(#1) }        % algebra of compact linear operators
\def\CL#1{\mathscr{C}(#1) }        % algebra of compact linear operators
\def\FO#1{\Phi(#1) }               % set of Fredholm operators
\def\FK#1#2{\Phi_{#2}(#1) }        % set of Fredholm operators with index k

\def\hilH{\mathit{H}}              % Hilbert space H
\def\clU{\mathit{U}}               % closed subspace U

\def\orthO{\mathfrak{O}}           % orthonormal system O
\def\orthS{\mathfrak{S}}           % orthonormal system S
\def\basB{\mathfrak{B}}            % basis B of a Hilbert space

\def\GL#1 {\mathbf{GL}(#1)}        % linear group
\def\OG#1 {\mathbf{O}(#1)}         % linear group
\def\SOG#1 {\mathbf{SO}(#1)}       % linear group

%------------------------------------------------------------------------------%
% distribution theory                                                          %
%------------------------------------------------------------------------------%
\def\disF{F}                  % distribution F
\def\disG{G}                  % distribution G
\def\disH{H}                  % distribution H
\def\disU{U}                  % distribution U
\def\disT{T}                  % distribution T
\def\SF{C^{\infty}(\G)}       % space of smooth functions
\def\TF{C_c^{\infty}(\G)}     % space of compactly supported smooth functions
\def\TFRd{C_c^{\infty}(\R^d)} % space of test functions
\def\SRd{\mathscr{S}(\R^d)}   % Schwartz space
\def\SR1{\mathscr{S}(\R)}     % Schwartz space
\def\B1#1{C^1(#1) }           % space of continuously differentiable functions
\def\Bm#1{C^m(#1) }           % space of continuously differentiable functions
\def\Ck{C^k(\G)}
\def\TD{\mathscr{S}'(\R^d)}   % space of temperated distributions
\def\E{\cal{E}}               % space of distributions with compact support
\def\D{\mathscr{D}}           % D
\def\FT{\mathcal{F}}          % Fourier transformation F
\def\FT{\mathscr{F}}          % Fourier transformation F

%------------------------------------------------------------------------------%
% other definitions                                                            %
%------------------------------------------------------------------------------%
\def\Proof{\textsc{Proof}}
\def\FRECHET{Fr\'{e}chet}
\def\ARZELA{Arzel\'{a}}
\def\GARDING{G\r{a}rding}
\def\HOELDER{H\"older}
\def\SCHROEDINGER{Schr\"odinger}
\def\JOERGENS{J\"orgens}
\def\WUEST{W\"uest}
\def\KAEHLER{K\"aehler}
\def\HOEFLINGER{H\"oflinger}
\def\POINCARE{Poincar\'{e}}
\def\FA{Functional Analysis}

\def\Author   {K.\,{\HOEFLINGER}}
\def\AuthorUC {K.\,HOEFLINGER}
\def\Version  {1.0}
\def\Email    {khoeflinger@t-online.de}


%\renewcommand\thesection{\Roman{section}}

\def\Title{Abstract Harmonic Analysis}

\titleformat{\section}
{\normalfont \bfseries \filcenter}
{\thesection.\;}{0mm}{}

\titleformat{\subsection}[runin]
{\normalfont \filcenter}
{}{0mm}{}

\titleformat{\subsubsection}[runin]
{\itshape \filcenter}
{}{0mm}{}

\titlespacing*{\section}      {5pt}{12pt}{ 4pt}
\titlespacing*{\subsection}   {0pt}{ 6pt}{12pt}
\titlespacing*{\subsubsection}{0pt}{ 6pt}{ 3pt}

%------------------------------------------------------------------------------%
%                                BEGIN OF DOCUMENT                             %
%------------------------------------------------------------------------------%
\begin{document}
\thispagestyle{empty}

\centerline{\large{\textbf{\Title}}}
\vspace{4pt}
\centerline{\small{{\textsc{\Author}}}}
\vspace{10pt}

\fancyhead[C] {\small{\textsc{\Title}}}
\fancyhead[OL]{\small{\thepage}}
\fancyhead[ER]{\small{\thepage}}

\myfootnote
{
  Version {\Version} from \today;
  Email:  \textit{khoeflinger@\,t-online.de}
}
The aim of this article
is to review some basic concepts of \textit{abstract harmonic analysis}.

\section{Basic Theory of Banach Algebras}
%------------------------------------------------------------------------------%
By a \textit{normed algebra} (rsp. \textit{Banach algebra}) is meant
a normed linear space (rsp. Banach space)
together with a multiplicative operation,
which is compliant with the norm in some sense.

\subparagraph{}
More specifically,
let $(\banB,+,\circ)$ be a $\F$-algebra,
that simultaneously constitutes a normed space $(\banB,\norm{\cdot})$.
Then $\banB$ is called a \textit{normed algebra},
if the following conditions are fulfilled:

\begin{itemize}[leftmargin=16mm,itemsep=2pt]
\item[(A-1)]
$\norm{x \circ y} \, \leq \, \norm{x} \cdot \norm{y}$
for each pair of elements $x,y \xeof \banB$.

\item[(A-2)]
If $\banB$ has an identity $\1 \xeof \banB$, then $\norm{\1} \eq 1$.
\end{itemize}

\subparagraph{}
Condition (A-1) implies
that the map $(x,y) \mapsto x \circ y$ is continuous.
If (A-2) is valid,
then $\banB$ is called a \textit{unital} algebra.
Finally,
$\banB$ is said to be \textit{commutative},
if
\begin{itemize}[leftmargin=16mm,itemsep=2pt]
\item[(A-3)]
\;$x \circ y \eq y \circ x$ \; for all $x,y \xeof \banB$.
\end{itemize}
As for normed linear spaces,
completeness of normed algebras is of interest,
thus a normed algebra is denoted a \textit{Banach algebra},
if it is complete as a normed linear space.

\subparagraph{} %- the unital Banach algebra of bounded linear operators
The most prominent example of a non-commutative Banach algebra is given by
the $\C$-algebra $\BO{\banX}$ of bounded linear operators
$\linA \:\banX \xto \banX$,
where $\banX$ is a real or complex Banach space.
It forms a unital Banach algebra {\wrt} composition of mappings
and the operator norm $\norm{\linA}$ of $\linA \xeof \BO{\banX}$.

\paragraph{} %- Gelfand-Mazur theorem
One of the key results in the theory of complex Banach algebras
is the \textit{Gelfand-Mazur theorem}
(see also \cite{VII_Ka}, p.\,14):

\begin{theorem1}\label{GELFAND_MAZUR_THEOREM}
If each element $x \neq 0$ in a complex Banach algebra $\banB$ is invertible,
then $\banB \cong \C$.
\end{theorem1}

\subparagraph{}
In other words, 
theorem \ref{GELFAND_MAZUR_THEOREM} quotes that
the field of complex numbers is essentially the only complex Banach algebra,
in which each element $x \neq 0$ is invertible.

\paragraph{} %- Neumann series
Let $\banB$ be a complex and unital Banach algebra.
By the \textit{Neumann series} of $x \xeof \banB$ is meant the infinite sum
$
s_x := \sum_{n=0}^{\infty} \, x^n.
$
If $s_x$ is convergent in $\banB$,
then the element $\1-x$ is invertible,
and a short computation proves $(\1-x)^{-1} \eq s_x$.
Like in the case $\banB \eq \C$,
the {Neumann series} is convergent for all $x \xeof \banB$
fulfilling $\norm{x} < 1$.

\subparagraph{}
If $x \xeof \banB$ is invertible,
then any other element $x' \xeof \banB$ having the property
\[
\norm{ x - x'} \, < \, \frac{1}{\snorm{x^{-1}}}
\]
proves to be invertible, too.
Thus all elements in a sufficiently small neighbourhood
of an invertible element have convergent Neumann series
and are invertible,
which implies that
the set $\GL{\banB} $ of invertible elements in $\banB$ is open.
It is called the \textit{group of units} in $\banB$,
since $x \circ y$ and $x^{-1}$ are invertible for $x,y \xeof \GL{\banB} $.
Both operations $(x,y) \mapsto x \circ y$
and $x \mapsto x^{-1}$ are continuous, 
so $\GL{\banB} $ forms a topological group,
when considered in the subset topology of $\banB$.

\paragraph{} %- examples of commutative Banach algebras
Here are some important examples of commutative Banach algebras
{\wrt} applications:

\begin{itemize}[leftmargin=12pt,itemsep=4pt]
\item[1.]
The field $\C$ of complex numbers may be regarded
as a complex unital Banach algebra in its own right.

\item[2.]
For any topological space $\topX$
the function spaces $C_b(\topX)$ and $C_0(\topX)$,
both endowed with the supremum norm $\norm{f} := \sup_{x} \abs{f(x)} $.

\subparagraph{}
If $\topX$ is \textit{compact},
then the function space $C(\topX)$ consisting of continuous functions
$f\: \topX \xto \C$ forms a Banach algebra,
which in addition is unital with identity $\1(x) := 1$ for $x \xeof \topX$.

\item[3.]
The algebra $L^1(\topG,\Borel{\topG},\mu)$ of Lebesgue-integrable functions
{\wrt} the Haar measure $\mu$ on a locally compact topological group $\topG$,
where multiplication is defined by \textit{convolution}
\[
(f \ast g)(x) \eqdef \int_{\R} f(y) \, g(x-y) \, d\mu
\]
for $f,g \xeof L^1(\topG)$.
Important examples for $\topG$ are the classical groups $\Rn$, $\Z$
and $\mathbb{T} \eq \{z \xeof \C: \norm{z} \eq 1\}$.

\item[4.]
The Banach algebra $L^{\infty}(X)$ of essentially bounded functions
$f \: X \xto \R$, 
where $(X,\Borel{X},\mu)$ may be any measure space.
Addition and multiplication in this algebra are defined pointwise.

\item[5.]
The algebra $\cal{M}^r(\topG)$ of regular complex Borel measures
defined on a locally compact topological group $\topG$.
The multiplication is defined by convolution of measures:
\[
(\mu \ast \nu)(B) \eqdef
(\mu \times \nu) (\{(x,y) \in \topG \xtimes \topG:x+y \in B\})
\]
for any Borel set $B \subset \topG$.
This Banach algebra is even unital, 
since the so-called \textit{Dirac measure} $\mu_0$ owns the property
$\mu_0 \ast \nu \eq \nu \ast \mu_0 \eq \nu$
for all $\nu \xeof \cal{M}^r(\topG)$.
\end{itemize}

\section{Gelfand Representation}
%------------------------------------------------------------------------------%
Given the Banach algebras $\banB_1, \banB_2$,
a linear mapping $\Phi\:\banB_1 \xto \banB_2$
is denoted a \textit{Banach algebra homomorphism},
if it preserves multiplication according to 
\[
\Phi(x \circ y) \= \Phi(x) \circ \Phi(y)
\fa x,y \in \banB_1.
\]
Moreover,
$\banB_1$ and $\banB_2$ are said
to be \textit{isomorphic} to each other, 
if there is a bijective homomorphism $\Phi \:\banB_1 \xto \banB_2$ fulfilling
$\norm{\Phi(x)}_{\banB_2} \eq \norm{x}_{\banB_1}$ for all $x \xeof \banB_1$,
which in this case is called a \textit{Banach algebra isomorphism}
(notation: $\banB_1 \cong \banB_2$).

\paragraph{} %- Gelfand representation
The \textit{Gelfand representation} is a powerful mapping
in the setting of commutative Banach algebras
and provides a Banach algebra homomorphism
between an arbitrary commutative Banach algebra $\banB$ and
the algebra $C_0(\topX)$ of continuous functions vanishing at infinity,
where $\topX$ is assumed to be a locally compact Hausdorff space.

\subparagraph{} %- characters on Banach algebras
A \textit{character} on a complex Banach algebra $\banB$
is a \textit{non-zero} Banach algebra homomorphism $\chi \:\banB \xto \C$.
Characters (often also called \textit{multiplicative functionals}) are
of great importance in the theory of unital commutative Banach algebras,
as we will see in the following.

\subparagraph{}
By the multiplication property of a Banach algebra homomorphism,
each character $\chi$ on $\banB$ satisfyies the inequality
$\norm{\chi} \leq 1$,
so $\chi$ is a continuous mapping and
thus an element of the dual space $\banB'$.
The set $\Gamma$ of all characters is called the \textit{structure space}
or the \textit{spectrum} of the Banach algebra $\banB$.

\subparagraph{} %- Gelfand topology
Due to the inequality $\norm{\chi} \leq 1$ for all $\chi \xeof \Gamma$,
the set $\Gamma$ is a subset of the closed unit ball $B_1(0) \subset \banB'$.
Thus we may endow $\Gamma$ with the subset topology on $B_1(0)$
inherited from the weak-*-topology on $\banB'$,
which is called the \textit{Gelfand topology} on the structure space.
Now,
by the \textit{Alaoglu-Bourbaki theorem}
the weak-*-topology on $\banB'$ is weak*-compact,
implying
that the Gelfand topology must be locally compact and Hausdorff.
Hence the structure space $\Gamma$ becomes into
a locally compact Hausdorff space.
Moreover, 
if $\banB$ is unital,
it can be shown that the set $\Gamma$ is even compact.

\paragraph{} %- Gelfand representation
From now on
let $\banB$ be a \textit{commutative} Banach algebra.
We consider the natural mapping $\gamma \:\banB \xto C(\Gamma)$,
where $\gamma$ assigns to each $x \xeof \banB$ the function
$\hat{x} \xeof C(\Gamma)$ by way of $\hat{x}(\chi) := \chi(x)$,
that is,
the function values of $\gamma$ are given by evaluating
the characters $\chi$ at the element $x$.
The mapping $\gamma$ is called the \textit{Gelfand representation}
of $\banB$.

\paragraph{} %- properties of the Gelfand representation
We collect the main properties of the Gelfand representation
for a commutative unital Banach algebra $\banB$ to the following list:

\begin{itemize}[leftmargin=12mm]
\item[1.]
$\gamma$ is a Banach algebra homomorphism
from $\banB$ to $C(\Gamma)$ or $C_0(\Gamma)$,
depending whether $\banB$ is unital or not.

\item[2.]
If $\banB$ is unital,
then $\hat{\1}$ is the constant function $1$ on $\Gamma$.

\item[3.]
If $\banB$ is unital,
then $x \xeof \banB$ is invertible
iff $\hat{x}$ nowhere vanishes.

\item[4.]
$\norm{ \hat{x} }_{\infty} \leq \norm{x}$ for each $x \xeof \banB$.
\end{itemize} 

\subparagraph{}
The Gelfand representation is injective
iff $\banB$ is a \textit{semisimple} Banach algebra,
that is,
if the intersection of all maximal ideals in $\banB$,
called the \textit{Jacobsen radical} of $\banB$,
is equal to the set $\{0\}$.

\section{Generalities on Topological Groups}
%------------------------------------------------------------------------------%
Before discussing abstract Fourier transformations,
it is useful to review some concepts coming from general topology,
especially that of a topological group, 
which plays a fundamental role in the theory of Fourier transforms.

\subparagraph{} %- topological groups
There is a subclass within that of general topological spaces,
which is of great importance in higher analysis
and whose members are called \textit{topological groups}.
They are mathematical objects owning the algebraic structure of a group
and a topological structure together.
Both have to be compatible to each other in the following sense:
let $(G,\circ)$ be an ordinary group and a topological space $(G,\topT)$
at once.
Then $G$ is called a \textit{topological group},
if the group operation $(x,y) \mapsto x \circ y^{-1}$ is continuous.

\subparagraph{}
Instead to require continuity of the mapping $(x,y) \mapsto x \circ y^{-1}$,
one can equivalently require
that the mappings $(x,y) \mapsto x \circ y$ and $x \mapsto x^{-1}$
have to be continuous.

\paragraph{} %- examples of topological groups
Well-known examples of topological groups are:

\begin{itemize}[leftmargin=12pt]
\item[1.]
each group $\grpG$ with its discrete topology $2^{\grpG}$,

\item[2.]
the group $\R^{\times} := \R - \{0\}$,

\item[3.]
the group $\C^{\times} := \C - \{(0,0)\}$,

\item[4.]
the group $\R^+ := \{x \xeof \R: x >0\}$,

\item[5.]
the additive groups $(\R,+)$ and $(\C,+)$,

\item[6.]
the \textit{circle group} $\mathbb{T} := \{z \xeof \C: \abs{z} \eq 1 \}$,

\item[7.]
each subgroup $\topU$ of a topological group $\topG$
endowed with the subspace topology
and its closure $\overline{\topU}$ are topological groups.
\end{itemize}

\paragraph{}
Like for purely algebraic groups,
a \textit{homomorphism} between topological groups
is defined to be a \textit{continuous} group homomorphism,
and such a homomorphism $\iota$ is called an \textit{isomorphism},
if it is a group isomorphism and $\iota^{-1}$ is continuous,
too
(we use the notation $\topG \cong \topH$ in this case).

\subparagraph{} %- (R,+) is isomorphic to (R+,.)
A well-known example of a group isomorphism is
the \textit{exponential function} $t \mapsto e^{t}$,
which maps the topological group $(\R,+)$ to the group $(\R^+,\cdot)$.
Rather similar,
the function
$t \xeof [0,1) \mapsto e^{2 \pi it}$
forms a group isomorphism between the topological groups
$\R / \Z$ and $\mathbb{T}$.

\paragraph{} %- automorphisms between topological groups
Because in a topological group $\topG$
the left  translation $x \mapsto g \circ x$ and
the right translation $x \mapsto x \circ g$
are continuous mappings,
they are both examples for \textit{automorphisms},
that is,
isomorphisms from the group $\topG$ to itself.
Hence the topological properties of a topological group
are completely determined by the topology in a neighbourhood
of the identity element $\1 \xeof \topG$,
which gives rise to call a topological group to be \textit{homogeneous}
in this case.

\subparagraph{}
Moreover,
the kind,
how the identity $\1$ of a topological group $\topG$ lies
in a given subgroup $\topU$,
determines the topology of that subgroup.
In particular,

\begin{itemize}[leftmargin=12pt]
\item[1.]
$\topU \subset \topG$ is open iff $\1$ is an inner point of $\topU$.

\item[2.]
$\topU$ is discrete iff $\1$ is an isolated point of $\topU$.
\end{itemize}

\paragraph{} %- circle group
Remember that the {circle group} $\mathbb{T}$ is the group
consisting of all complex numbers of modulus $1$ with group operation
\[
e^{i \alpha} \circ e^{i \beta} \eqdef
e^{i \alpha} \cdot e^{i \beta} \eqdef
e^{i(\alpha + \beta)},
\quad \alpha,\beta \xeof \R.
\]
Endowed with the subset topology inherited
from the standard topology on the field $\C$,
the unit circle $\mathbb{T}$ forms a topological group.
Consider the set $C(\topG,\mathbb{T})$ consisting of all
continuous functions $f \:\topG \xto \mathbb{T}$.
By furnishing this set with the compact-open topology,
it becomes into a Hausdorff space.

\subparagraph{} %- characters in the setting of topological groups
A \textit{character} $\chi$ is a (topological) group homomorphism from a
topological group $\topG$ to the circle group $\mathbb{T}$
and hence a continuous mapping.
We denote $\topG^*$ the set of characters $\chi \:\topG \xto \mathbb{T}$,
which forms a subset of the Hausdorff space $C(\topG,\mathbb{T})$.
We can make $\topG^*$ into a topological space by taking the subset topology
{\wrt} the compact-open topology on $C(\topG,\mathbb{T})$.

\paragraph{} %- dual groups
Performing multiplication of characters $\chi_1,\chi_2 \xeof \topG^*$ by
\[
\chi_1 \circ \chi_2\: g \mapsto \chi_1(g) \cdot \chi_2(g),
\]
we obtain a further character $\chi_1 \circ \chi_2 \xeof \topG^*$.
It can be shown that this operation has all the properties of a group,
and since it is commutative,
$\topG^*$ has the structure of an Abelian group,
called the \textit{character group} or the \textit{dual group} of $\topG$.
Moreover,
the group operation in $\topG^*$ is continuous,
hence the dual group of any topological group forms a topological group
as well.

\paragraph{}
In the following table we list three classical topological Abelian groups
together with their dual groups and characters:

\vspace{3mm}
\renewcommand{\arraystretch}{1.2}\normalsize
\centerline
{
\begin{tabular}{|c|c|c|}
\hline
$\topG$ & $\topG^*$       & characters                           \\
\hline
$\R$    & ${\R}^* \cong \R$ & $\chi_x(t) := e^{itx}, \, x \in \R$  \\
$\Z$    & ${\Z}^* \cong \mathbb{T}$ & $\chi_z(n) := e^{inz},
                 z \in \mathbb{T}$                                \\
$\mathbb{T}$    & ${\mathbb{T}}^* \cong \Z$ & $\chi_n(z) := z^n, 
                 \, n \in \Z$                                     \\
\hline
\end{tabular}
}

\vspace{3mm}
\centerline{\small Table 1: classical topological groups and dual groups.}
\vspace{2mm}

\paragraph{}
Concerning cartesian products of topological groups,
for each topological group $\topG$ and $\topH$
we have
\[
(\topG \times \topH)^*
\, \cong \,
\topG^* \times \topH^*.
\]
As a consequence,
the group isomorphisms ${\Rn}^* \cong \Rn$ hold for any integer $n$.

\section{Locally Compact Abelian Groups}
%------------------------------------------------------------------------------%
Recall that a topological space is called \textit{locally compact},
if every subset owns a compact neighbourhood.
The abbreviation \textit{LCA group} usually stands
for a \textit{locally compact Abelian group}.
The meaning of locally compact groups is based on the fact
that these groups have non-trivial characters,
that is,
we have $\topG^* \neq \{\chi_1\}$ for such a group ${\topG}$,
where $\chi_1$ is the character with $\chi_1(x) \eq 1$ for all $x \xeof \topG$.

\subparagraph{}
Suppose that the Abelian group ${\topG}$ is \textit{locally compact}.
Then it can be shown (for example by the {\ARZELA}-Ascoli theorem)
that the dual group $\topG^*$ is also locally compact {\wrt}
its compact-open topology.
Let $\topG^{**}$ denote the \textit{bidual group} 
of a topological group $\topG$,
which is defined to be the dual group of its character group.

\paragraph{} %- Pontryagin's duality theorem
One of the main results in the theory of LCA groups is the famous
\textit{Pontryagin duality theorem}:

\begin{theorem1}\label{PONTRYAGIN_DUALITY_THEOREM}
If $\topG$ is LCA,
then the groups $\topG$ and $\topG^{**}$ are isomorphic to each other.
\end{theorem1}

\subparagraph{}
Notice that
each group element $g \xeof \topG$ gives rise to a character $\chi_x$
on the dual group $\topG^*$ by
$\chi_x(\xi) := \xi(x)$ for all $\xi \xeof \topG^*$.
Thus the mapping $x \mapsto \chi_x \xeof \topG^{**}$ is an injective group
homomorphism,
and Pontryagin duality means
that the canonical embedding $\iota \: \topG \xto \topG^{**}$,
$x \mapsto \chi_x$ is also surjective.

\subparagraph{}
It is easy to see
that the groups listed in table 1 are examples of Pontryagin duality.
This kind of duality also implies remarkable relationships between
the topological properties of a LCA group $\topG$ and its dual group $\topG^*$:

\begin{itemize}[leftmargin=15pt]
\item[1.]
$\topG$ is compact \iff $\topG^*$ is discrete.

\item[2.]
$\topG$ is discrete \iff $\topG^*$ is compact.
\end{itemize}

\paragraph{}
Finally,
let us mention
that it is possible to endow the convolution algebra $L^1(\topG,+,\ast)$
of any LCA group $\topG$
with an involutional mapping $* \:L^1(\topG) \xto L^1(\topG)$,
defined for $f \xeof L^1(\topG)$ by
\[
f^*(x) \eqdef \overline{f(-x)}, \quad \quad (x \in \topG).
\]

\section{Abstract Fourier Transformation}
%------------------------------------------------------------------------------%
Since $L^1(\topG,+,\ast)$ forms a commutative Banach algebra for every
LCA group $\topG$,
we may study its Gelfand transformation in detail,
which in this case is denoted \textit{Fourier transformation}.
If $\Gamma_{\topG}$ is the structure space of the algebra $L^1(\topG,+,\ast)$,
then Fourier transformation maps each function $f \xeof L^1(\topG)$
to a function $\hat{f} \xeof C_0(\Gamma_{\topG})$.
It constitutes a Banach algebra homomorphism
between $L^1(\topG)$ and $C_0(\Gamma_{\topG})$,
where convolution in $L^1(\topG)$ corresponds to multiplication
in $C_0(\Gamma_{\topG})$.

\paragraph{}
Moreover,
each character $\chi \xeof \topG^*$ gives rise to
a continuous multiplicative functional $\gamma \xeof \Gamma_{\topG}$,
that is defined according to
\[
\gamma(\chi) \eqdef \int_{\topG} \, f \, \overline{\chi} \, d\mu,
\]
and reversely each continuous multiplicative functional
$\gamma \xeof \Gamma_{\topG}$ can be represented in this way,
Thus,
we identify the structure space of the convolution algebra $L^1(\topG,+,\ast)$
with the dual group $\topG^*$,
and Fourier transformation appears as the Gelfand transformation
up to this identification:

\begin{itemize}[leftmargin=12pt]
\item[1.]
For $f \xeof L^1(\topG)$
the mapping $\hat{f} \: \widehat{\topG} \mapsto \C$ according to
\[
\hat{f}(\chi) := \int_{\topG} \, f \, \overline{\chi} \, d\mu,
\quad
\chi \in \widehat{\topG}
\]
is called the \textit{Fourier transform} of $f$.

\item[2.]
The mapping
$\FT \: f \xeof L^1(\topG) \xto \hat{f} \xeof C_0(\topG^*)$
is called \textit{Fourier transformation}.
\end{itemize}

\paragraph{} %- properties of Fourier transformation
Let us repeat some properties of Fourier transformation
that are already known from the general Gelfand transformation:

\begin{itemize}[leftmargin=12pt]
\item[1.]
For every $f \xeof L^1(\topG)$ the mapping $\hat{f}$ is continuous on $\topG^*$.

\item[2.]
the function $\hat{f}$ is vanishing at infinity
(\textit{Riemann-Lebesgue theorem}),
that is,
it maps absolutely integrable functions
to elements of the function space $C_0(\topG^*)$.

\item[3.]
$\FT$ is injective,
that is,
$\FT(f) \neq \FT(g)$ for $f \neq g$.

\item[4.]
$\FT$ is norm-preserving,
that is,
$\norm{\FT(f)}_{\infty} \leq \norm{f}_1$.

\item[5.]
$\FT(\alpha f + \beta g) \eq \alpha \FT(f) + \beta \FT(g)$,
    $\alpha,\beta \xeof \R$.

\item[6.]
$\FT(f \ast g) \eq \sqrt{2\pi} \, \FT(f) \cdot \FT(g)$.
\end{itemize}

\subparagraph{}
In addition,
the image of $L^1(\topG)$ under $\FT$ forms a separating subalgebra
of $C_(\Gamma_{\topG})$,
which lies even dense in $C_0(\Gamma_{\topG})$. 

\paragraph{} %- Fourier transformation for the groups R,Z,T
It is nearby to consider Fourier transformation
on $L^1(\topG,+,\ast)$ for the classical groups
$\R$, $\R^n$, $\Z$ and $\mathbb{T}$:

\begin{itemize}[leftmargin=12pt]
\item[1.]
If $\topG \eq \R$,
then the character group of $(\R,+)$ is the group $(\R,+)$ itself.
Hence a function $f \xeof L^1(\R)$ owns the Fourier transform
$\hat{f} \:\R \xto \C$,
whose function values are defined as follows:
\[
\hat{f}(x) \eqdef
\frac{1}{\sqrt{2\pi}} \int_{- \infty}^{\infty} f(t) \, e^{-itx} \, dt,
\quad (x \xeof \R).
\]

\item[2.]
In case of $\topG \eq \R^n$
the character group of $(\R^n,+)$ coincides with the group $(\R^n,+)$ as well.
The Fourier transform of $f \xeof L^1(\R^n)$ is the function defined by
\[
\hat{f}(x) \eqdef
\frac{1}{(2\pi)^{n/2}}
\int_{- \infty}^{\infty} f(t) \, e^{-i \scalar{t}{x}} \, dt,
\quad (x \xeof \R^n).
\]

\item[3.]
Let $\topG \eq \Z$.
The normed Haar measure on this group is nothing but the counting measure,
and the character group $(\Z,+)$ is isomorphic to the multiplicative group
$(\mathbb{T},\cdot)$.
Hence the Fourier transform of a sequence $\{f_n\}_{n \in \N}$
is the mapping $\hat{f} \: \mathbb{T} \xto \C$,
which is defined pointwise by the \textit{Fourier series}
\[
\hat{f}(\omega) \eqdef
 \frac{1}{\sqrt{2\pi}} \sum_{n= - \infty}^{\infty} f_n \, e^{-i n \omega},
 \quad (\omega \in \mathbb{T}).
\]

\item[4.]
The character group of the group $(\mathbb{T},\cdot)$ is isomorphic
to the group $(\Z,+)$,
hence the Fourier transform of a function
$f \: \mathbb{T} \xto \R$ is the sequence $\hat{f}:\Z \xto \C$ with members
\[
\hat{f}_n \eqdef
\frac{1}{\sqrt{2\pi}} \int_{\mathbb{T}} f(z) \, e^{-inz} \, dz,
\quad (n \in \Z).
\]
\end{itemize}

\paragraph{} %- inverse Fourier transformation
By \textit{inverse Fourier transformation} is meant the procedure,
how certain functions in the algebra $L^1(\topG)$ can be retrieved
from the function values of the corresponding Fourier transforms.

\paragraph{}
Since the adjoint group $\topG^*$ of a LCA group $\topG$
with Haar measure $\mu$ is LCA again,
it owns a Haar measure as well,
hence the convolution algebra $L^1(\topG^*)$ is well-defined.
Let $f$ be an element of $L^1(\topG)$
fulfilling the additional condition $\hat{f} \xeof L^1(\topG^*)$.
Then there exists the Fourier transform of $\hat{f}$,
\[
\hat{\hat{f}}(x) \= 
\int_{\topG^*} \hat{f}(\gamma) \overline{x(\gamma)} \, d\nu,
\]
which forms a complex-valued mapping
defined on the biadjoint group $\topG^{**}$.
Due to Pontryagin duality,
$\topG^{**}$ is isomorphic to $\topG$,
and we may define the \textit{inverse Fourier transform} $F$
according to
\[
F(x) \= \int_{\topG^*} \hat{f}(\gamma) \, \overline{\gamma(x)} \, d \mu,
\quad (x \in \topG).
\]
It can be shown that the functions $f$ and $F$ are $\mu$-nearly everywhere
identical on $\topG$.

\subparagraph{}
Especially for the group $(\R,+)$,
where even $\R^*$ may be identified with $\R$,
we obtain the inverse Fourier transform of $\hat{f} \xeof L^1(\R)$ according to
\[
F(x) \= \int_R \hat{f}(t) \, e^{itx} \, dt, \quad (x \in \R).
\]

\paragraph{} %- the algebra of finite Borel measures
Let $\topX$ be any topological space and $\mu$ be a measure
define on the Borel $\sigma$-algebra $\Borel{\topX}$.
A Borel set $B \subseteq \topX$ is called \textit{inner regular} if
\[
\mu(B) \= \sup \,
          \{\mu(\compK): \compK \subset B, \, \compK \mbox{ kompact}\}.
\]
Then a measure $\mu$ is denoted as \textit{inner regular},
if all Borel sets are so.
Hence inner regular measures defined on $\topX$ are already uniquely
defined by their values for compact subsets of $\topX$.
Each inner regular measure on a topological space is said to be \textit{Radon}
if it adopts finite values for compact sets.

\subparagraph{}
Rather similar,
a measure $\mu$ is called \textit{outer regular},
if for each Borel set $B \subseteq \topX$ the equality
\[
\mu(B) \= \inf \,
          \{\mu(\openO): B \subset \openO, \, \openO \mbox{ open} \}.
\]
holds.
Finally,
a measure is called \textit{regular},
if it is inner and outer regular at the same time.

\paragraph{} %- convolution algebras of measures
Let $\MCR{\topG}$ be the Banach space of finite regular Borel measures,
where $\topG$ is a LCA group.
Let $\mu_1, \mu_2$ be measures in $\MCR{\topG}$ and let $\mu_1 \xtimes \mu_2$
be the product measure defined on the group $\topG \xtimes \topG$.
If $B \xeof \Borel{\topG}$ is measurable,
then the set $B'$ is defined by
\[
B' \eqdef \{(x,y) \in \topG \times \topG: x+y \in B\}.
\]
This set is measurable {\wrt} the $\sigma$-Algebra
in $\topG \xtimes \topG$.
Moreover,
the \textit{convolution} operation $\mu_1 * \mu_2$
of two measures $\mu_1$ und $\mu_2$ is defined to be the measure
\[
(\mu_1 * \mu_2)(B) \eqdef (\mu_1 \times \mu_2)(B')
\quad (B \in \Borel{\topG}),
\]
which owns the following properties:

\begin{itemize}[leftmargin=12pt]
\item[1.]
$\mu_1 * \mu_2 \xeof \MCR{\topG}$,

\item[2.]
$(\mu_1,\mu_2) \mapsto \mu_1 * \mu_2$ is commutative and associative,

\item[3.]
the inequality
$\norm{\mu_1 * \mu_2} \, \leq \, \norm{\mu_1} \cdot \norm{\mu_2}$ holds.
\end{itemize}

\subparagraph{} %- Dirac measure
Hence the Banach algebra $\MCR{\topG,+,*}$ is unital and commutative.
Its unit element is the \textit{Dirac measure} $\delta_{\0}$,
which distinguishes from other measures in this algebra by the property
that it is concentrated to the element $\0 \xeof \topG$.

\paragraph{}
In case of $\topG=(\Rn,+)$,
we may regard the algebra $\MCR{\Rn,+,*}$ as the Banach algebra,
that is norm-isomorphic to the completion
of the non-unital convolution algebra $L^1(\Rn,+,*)$.

\paragraph{} %- Fourier-Stieltjes transformation
\textit{Fourier-Stieltjes transformation} is another phrase for
Gelfand transformation on the unital and commutative Banach algebra
$\MCR{\topG,+,*}$ of finite and regular complex Borel measures
on a LCA group $\topG$.
In fact,
if $\mu \xeof \MCR{\topG,+,*}$,
then the mapping $\hat{\mu}$ defined on the dual group $\topG^*$ by
\[
\hat{\mu}(\gamma) \eqdef
\int_{\topG} -\gamma(x) \, d\, \mu(x), \quad \quad \gamma \in \topG^*
\]
is called \textit{Fourier-Stieltjes transform} of  $\mu$.

\paragraph{} %- properties of Fourier-Stieltjes transformation:
Let $B(\topG^*)$ be the set of functions $\hat{\mu}$
with $\mu \xeof \MCR{\topG,+,*}$.
Then the following properties are valid:

\begin{itemize}[leftmargin=12pt]
\item[1.]
Each function $\hat{\mu} \xeof B(\topG^*)$ is bounded and uniformly continuous.

\item[2.]
If $\sigma \eq \mu * \nu$,
then $\hat{\sigma} \eq \hat{\mu} \cdot \hat{\nu}$ holds,
that is,
the mapping $\mu \mapsto \hat{\mu}(\gamma)$ 
is a Banach algebra endomorphism on $\MCR{\topG,+,*}$
for every $\gamma \xeof \topX^*$.

\item[3.]
The set $B(\topG^*)$ is invariant under translation, multiplication and
complex conjugation.
\end{itemize}

\section{Fourier Transformation in Euclidean Spaces}
%------------------------------------------------------------------------------%
To the end we elaborate \textit{Fourier transformation} in the spaces $\Rd$,
which is an important tool in advanced analysis with many useful applications,
especially to the study of
linear differential operators with constant coefficients.

\paragraph{} %- convolutions
Rather Fourier transformation can be defined in various linear spaces,
it is customary to introduce this mapping first in Lebesgue space $L^1(\Rd)$.
Recall that this space contains all measurable functions $f \: \R \xto \C$
having the property
\[
\int_{\,\Rd} \, \abs{f} \, dx \, < \, \infty.
\]

\subparagraph{} %- convolution in L^1
Preliminary,
for elements $f,g \xeof L^1(\Rd)$ 
a multiplication $f \ast g$ is defined by
\[
(f \ast g) \, (x) \eqdef \int_{\,\Rd} f(t) \cdot g(x-y) \, dy
\]
and denoted the \textit{convolution} in $L^1(\Rd)$.
This mapping is commutative, associative and fulfills
for all $f,g \xeof L^1(\Rd)$ the estimate
\[
\norm{f \ast g}_1 \, \leq \, \norm{f} \cdot \norm{g}.
\]
Since $L^1(\Rd)$ is complete,
the algebra $(L^1(\Rd),+,\ast)$
constitutes a commutative but non-unital Banach algebra.

\paragraph{} %- Fourier transformation in L1
L1-Fourier transformation in Euclidean space $\Rd$
is nothing but the Gelfand mapping
$\FT \: f \xto \hat{f}$ \xeof $C_0(\Rd)$,
which sends each function $f \xeof L^1(\Rd)$
to the function $\hat{f} \xeof C_0(\Rd)$,
called the \textit{Fourier transform} of $f$ and defined by
\begin{equation}\label{VII_L1_FOURIER_TRANSFORM}
\hat{f}(\xi) \eqdef
\frac{1}{(2 \pi)^{d/2}}
\int_{\,\Rd} \, e^{-2 \pi i \scalar{\xi}{x}} \, f(x) \, dx
\quad \quad (\xi \xeof \Rd),
\end{equation}
where $\scalar{\xi}{x}$ is the usual inner product of $\xi,x \xeof \Rd$.

\paragraph{} %- properties of Fourier transformation in L1
Let us again collect the remarkable properties of the mapping $\FT$
in the special case $\topG \eq \topG^* \eq \Rd$:

\begin{itemize}[leftmargin=12pt,itemsep=5pt]
\item[1.]
For each $f \xeof L^1(\Rd)$ the function $\hat{f}$ is continuous on $\Rd$,
so $\hat{f} \xeof C(\Rd)$.

\item[2.]
$\hat{f}$ is vanishing at infinity,
so $\hat{f} \xeof C_0(\Rd)$.

\item[3.]
$\FT$ is injective,
that is,
$\FT(f) \neq \FT(g)$ for $f \neq g$.

\item[4.]
$\norm{\FT(f)}_{\infty} \leq \norm{f}_1$ for $f \xeof L^1(\Rd)$.

\item[5.]
$\FT(\alpha f + \beta g) \eq \alpha \, \FT(f) + \beta \, \FT(g)$
for $\alpha,\beta \xeof \C$,
thus $\FT$ is linear.

\item[6.]
$\FT(f \ast g) \= \FT(f) \cdot \FT(g)$,
so $\FT$ is an algebra homomorphism between the Banach algebras
$L^1(\Rd,+,\ast)$ and $C_0(\Rd,+,\cdot)$.

\item[7.]
If $f \xeof L^1(\Rd)$ and
$g \: x \xeof \Rd \mapsto \abs{x} f(x) \in L^1(\Rd)$,
then $\hat{f} \xeof C^1(\Rd)$ and
\[
\partial_j f(x) \= \widehat{-i \, \xi_j f(\xi)}
\quad (j=1,\ldots,d).
\]
\end{itemize}

\paragraph{} %-- Fourier-Plancherel transformation
%------------------------------------------------------------------------------%
We now introduce Fourier transformation in the space $L^2(\Rd)$.
Indeed,
Recall that
this space consists of the measurable functions $u \: \Rd \xto \C$ fulfilling
\[
\int_{\,\Rd} \, \abs{u} ^2 \, dx \, < \, \infty,
\]
and the inner product in $L^2(\Rd)$ is given by
\[
\scalar{u}{v} \= \int_{\,\Rd} \, u\,\overline{v} \, dx.
\]
If the functions $u,v \: \Rd \xto \C$ are Lebesgue-integrable
then by definition they belong to the space $L^1(\Rd)$,
and their \textit{Fourier transforms} are well-defined
by formula \eqref{VII_L1_FOURIER_TRANSFORM}.

\subparagraph{}
However, 
if is not only a member of $L^1(\Rd)$, 
but also $u \xeof L^2(\Rd)$,
we define Fourier transformation a little bit different:
\begin{equation}\label{VII_L2_FOURIER_TRANSFORM}
\hat{u}(\xi) \eqdef
\frac{1}{(2 \pi)^{d/2}}
\int_{\,\Rd} \, e^{-2 \pi i \scalar{\xi}{x}} \, u(x) \, dx.
\end{equation}
Then for functions $u,v \xeof L^2(\Rd)$ 
the inner product between $u$ and $v$
and the inner product of their Fourier transforms $\hat{u}$ and $\hat{v}$
are related by \textsc{Plancherel}\textit{'s formula}
\[
\scalar{u}{v} \= \scalar{ \hat{u} }{\hat{v} }.
\]
Hence \textit{Fourier transformation} $\FT \: u \mapsto \hat{u} $
constitutes a partial isometry from the space
$L^{1,2}(\Rd) := L^1(\Rd) \, \cap \, L^2(\Rd)$
to the space $L^2(\Rd)$.

\subparagraph{}
It is a crucial fact now
that $L^{1,2}(\Rd)$ proves to be dense in $L^2(\Rd)$.
Since the mapping $\FT$ is linear and continuous on $L^{1,2}(\Rd)$ ,
it is even \textit{uniformly continuous} on that space
and can be extended to an \textit{unitary},
that is,
a bijective and isometric linear operator
$\FT_P \: L^2(\Rd) \xto L^2(\Rd)$,
called the \textsc{Fourier-Plancherel} \textit{operator}
or the \textit{$L^2$-Fourier transformation} in $L^2(\Rd)$.

\subparagraph{}
But in general a function $u \xeof L^2(\Rd)$ is not an element of $L^1(\Rd)$.
Thus it is not possible to define the $L^2$-Fourier transform $\hat{u}$
of every $u \xeof L^2(\Rd)$
by means of formula \eqref{VII_L2_FOURIER_TRANSFORM}.
However,
for $u \xeof L^2(\Rd)$ and $R > 0$ the function
\[
\hat{u} _R(\xi) \eqdef
\frac{1}{(2 \pi)^{d/2}}
\int_{\,\abs{x} \leq R}
e^{-2 \pi i \scalar{x}{\xi}} \, u(x) \, dx
\]
is well-defined,
and one may obtain the $L^2$-Fourier transform $\hat{u}$ 
as the limit of the functions $\hat{u} _R$ for $R \xto \infty$
{\wrt} the $2$-norm,
that is,
$\lim_{\,R \to \infty} \, \norm{ \hat{u} - \hat{u} _R}_2 \= 0$.

\paragraph{} %- Bessel potential spaces
The Fourier-Plancherel operator can be used
to introduce the class of \textit{Bessel potential spaces}.
First recall
that the mapping $\FT_P$ is well-defined in Hilbert space $L^2(\Rd)$,
therefore it has also a meaning in the linear subspaces $H^k(\Rd)$
for $k \xeof \N$.
We define for each $k \xeof \N$ the real-valued mapping
\begin{equation}\label{VII_DEFINITION_OF_FUNCTION_P}
p_k:x \in \Rd \mapsto (1 + \abs{x} ^{2k})^{1/2}
\end{equation}
and the space
\[
\cal{H}^k(\Rd) \eqdef
\{ f \in L^2(\Rd): p_k \, \hat{f} \in L^2(\Rd) \},
\]
which is Hilbertian with inner product
$\scalar{f}{g}_{p_k} := \scalar{p_k \hat{f}}{p_k \hat{g}}$
for $f,g \xeof \cal{H}^k(\Rd)$ and derived norm
\[
\norm{f}_{p_k}^2 \eqdef \scalar{f}{f}_{p_k} \=
\int_{\,\Rd} \, (1+\abs{x} ^{2k}) \, \abs{\hat{f}} ^2 \, dx.
\]

\subparagraph{}
It turns out
that the so-called \textit{Bessel potential space}
$\cal{H}^k(\Rd)$ is the range of
the \textsc{Fourier-Plancherel} operator $\FT_P$
when restricted to $H^k(\Rd) \subset L^2(\Rd)$.
Hence $\FT_P$ is a bijective and norm-preserving mapping
from $H^k(\Rd)$ to $\cal{H}^k(\Rd)$, {\ie}
\[
\abs{f} _k \= \lVert \hat{f} \rVert _{p_k}
\quad \, \qfa f \in H^k(\Rd).
\]

\subparagraph{}
In order to extend the definition of the spaces $H^k(\Rd)$
to indices $s \xeof \R^+$,
let $p_s$ be defined as in \eqref{VII_DEFINITION_OF_FUNCTION_P}
but with $s > 0$ instead of $k \xeof \N$.
We define
\[
\cal{H}^s(\Rd) \eqdef
\{ \, f \in L^2(\Rd): p_s \hat{f} \in L^2(\Rd) \, \},
\]
then $\cal{H}^s(\Rd)$ given this way
coincides for $s \eq k \xeof \N$ with the classical Sobolev space $H^k(\Rd)$,
which motivates to define $H^s(\Rd) := \cal{H}^s(\Rd)$ for $s \notin \N$.
For $s<0$,
however,
we introduce $H^s(\Rd) := H^{-s}(\Rd)^*$.

\paragraph{} %-- Fourier transformation in the space of tempered distributions
We finally elaborate that Fourier transformation can be extended
to the space $\TD$ of tempered distributions.

\subparagraph{}
First notice 
that the Fourier transform $\hat{\phi}$
of a Schwartz function $\phi \xeof \SRd$ is well defined,
since $\SRd$ is a linear subspace of the convolution algebra $L^1(\Rd,+,\ast)$.
Moreover, 
Fourier transformation restricted to the Schwartz space $\SRd$
forms a linear mapping $\FT \:\SRd \xto \SRd$,
which sends each function $\phi \xeof \SRd$ to a further function
$\hat{\phi} \xeof \SRd$ according to the well-known formula
\begin{equation}\label{VII_FOURIER_TRANSFORM_OF_SCHWARTZ_FUNCTION}
\hat{\phi}(x) \= \FT(\phi)(x) :=
\frac{1}{(2\pi)^{d/2}}
\int_{\,\Rd} e^{-i <\xi,x>} \phi(\xi) \, d \xi
\quad (x \xeof \Rd).
\end{equation}

\subparagraph{}
The extraordinary meaning of Schwartz functions is founded by the fact that
the Fourier transformation has excellent properties on the space $\SRd$.
More precisely,
besides the well-known features of Fourier transformation in $L^1(\Rd,+,\ast)$,
the mapping $\FT$ is \textit{bijective} on $\SRd$.
Hence the \textit{inverse Fourier transform} exists for all $\phi \xeof \SRd$
and is described by the \textit{Fourier inversion formula}
\[
\check{\phi}(\xi) \= 
\FT^{-1}(\phi)(\xi) \eqdef
\frac{1}{(2\pi)^{d/2}}
\int_{\,\Rd} e^{i \scalar{x}{\xi}} \, \phi(x) \, dx
\quad \quad (\xi \xeof \Rd).
\]

\subparagraph{}
We next extend Fourier transformation in $\SRd$ to the space $\TD$
of tempered distributions by transposition,
that is
\[
(\FT \disT) \, \phi \, := \, \disT(\FT \phi),
\quad \quad \; \phi \xeof \SRd, \; \disT \xeof \TD.
\]
This mapping is bijective in the space $\TD$
and even continuous {\wrt} the weak*-topology on $\TD$.
The inverse Fourier transform of a tempered distribution $\disT$
is represented by the formula
\[
\FT^{-1}(\disT)(\phi) \, := \, \disT(\FT^{-1} \phi)
\quad \quad (\phi \xeof \SRd).
\]
For example,
the Fourier transform $\FT(\delta_0)$ of the Dirac distribution
$\delta_0 \xeof \cal{S}'(\R)$ is the continuous linear functional
\[
\FT(\delta_0)(\phi) \, := \,
\delta_0(\FT(\phi)) \= \hat{\phi}(0) \=
\frac{1}{\sqrt{2\pi}} \, \int_{\,\R} \phi(x) \, dx,
\quad \phi \xeof \SR1\,,
\]
so $\FT(\delta_0)$ is nothing but the distribution induced by
the constant function $x \mapsto \frac{1}{\sqrt{2\pi}}$.

% BIBLIOGRAPHY
%------------------------------------------------------------------------------%
\vspace{5mm}
\begin{thebibliography}{99}
\begin{small}

\bibitem{VII_Ka} E.\,Kaniuth:
\textit{A Course in Commutative Banach Algebras},
Springer GTM 246, 2009.

\bibitem{VII_Ru} W.\,Rudin:
\textit{Fourier Analysis on Groups},
Wiley-Interscience, 1962.

\bibitem{VII_SW} E.\,Stein, G.\,Weiss:
\textit{Introduction to Fourier Analysis on Euclidean Spaces},
Princeton Mathematical Series 32, Princeton University Press, 1971.

\end{small}
\end{thebibliography}

%------------------------------------------------------------------------------%
%                               END OF DOCUMENT                                %
%------------------------------------------------------------------------------%
\end{document}
